\documentclass[../../../include/open-logic-section]{subfiles}

\begin{document}

\olfileid{sth}{story}{rus}
\olsection{Paradoks Russell (maneh)}

Ing \olref[sfr][][]{part}, kita nggunakake teori himpunan naif.
Nanging miturut konsepsi kang \emph{luwih naif maneh}, himpunan
iku mung ekstensi saka predikat. Pamawas naif iki nuntut
prinsip ing ngisor iki:

\begin{defish}
  \emph{Komprehensi Naif.} $\Setabs{x}{\phi(x)}$ ana kanggo
  saben formula $\phi$.
\end{defish}

Sanajan prinsip iki katon narik kawigaten, bisa dibuktekake
yen prinsip iki ora konsisten. Kita wis weruh iki ing
\olref[sfr][set][rus]{sec}, nanging asil iki wigati banget
lan prasaja banget nganti pantes dibaleni, persis kaya sadurunge.

\begin{thm}[Paradoks Russell]
Ora ana himpunan $R = \Setabs{x}{x \notin x}$.
\end{thm}

\begin{proof}
Yen $R = \Setabs{x}{x \notin x}$ ana, mula $R \in R$ yen lan
mung yen $R \notin R$, sawijining kontradiksi.
\end{proof}

Russell nemokake asil iki ing Juni 1901. (Nanging dheweke
ora ngandharake paradoks ing wujud kang persis kaya kang
anyar wae kita aturake, amarga dheweke nimbang teori himpunan
Frege kaya kang diandharake ing \emph{Grundgesetze}. Kita bakal
bali menyang iki ing \olref[blv]{sec}.) Russell nulis surat
marang Frege tanggal 16 Juni 1902, nerangake inkonsistensi
sisteme Frege. Kanggo korespondensi kasebut lan sawetara
katrangan latar, delengen \citet[pp.~124--8]{Heijenoort1967}.

Perlu ditekanake yen bukti rong larik iki minangka asil saka
\emph{logika murni}. Pancen, ing notasi
$\Setabs{x}{x \notin x}$ kita sacara implisit nggunakake
aksioma Ekstensionalitas (apa aksioma iki dudu logis?).
Amarga $\Setabs{x}{\phi(x)}$ dimaksudake dadi himpunan
\emph{tunggal} (miturut Ekstensionalitas) kang ngemot
persis $\phi$s, yen himpunan iku ana. Nanging kita bisa
nyingkiri malah kesan yen Ekstensionalitas dibutuhake,
kanthi ngandharake asil iki kaya mangkene: \emph{ora ana
himpunan kang anggotane persis himpunan kang ora dadi
anggota awake dhewe}. Bab iki ora gumantung banget marang
himpunan. Kaya kang dingerteni Russell dhewe, penalaran kang
padha persis uga menehi konklusi: \emph{ora ana priya kang
nyukur persis para priya kang ora nyukur awake dhewe}.
Utawa: \emph{ora ana asu pug kang ngambu persis asu-asu pug
kang ora ngambu awake dhewe}. Lan sapiturute. Kanthi skema,
wujud asile mung:
\[
\lnot \exists x \forall z(Rzx \liff \lnot R zz).
\]
Lan iku sawijining teorema (skema) logika tataran kapisan.
Mula kita ora bisa nyingkiri Paradoks Russell mung kanthi
ngowahi teori himpunan: paradoks iki tuwuh sadurunge kita
malah \emph{mlebu} ing teori himpunan. Yen kita arep nggunakake
logika tataran kapisan (klasik), kita kudu \emph{nampa}
yen ora ana himpunan $R = \Setabs{x}{x\notin x}$.

Akibate mangkene. Yen kita arep nampa Komprehensi Naif
nanging \emph{nyingkiri} inkonsistensi, kita ora cukup mung
ngowahi \emph{teori himpunan}. Nanging kita kudu ngowahi
\emph{logika} kang digunakake kanthi dhasar.

Mesthi wae, teori himpunan kanthi logika nonklasik wis ana.
Nanging teori-teori iku---paling ora---dudu pendekatan baku.
Pendekatan baku tumrap Paradoks Russell yaiku nganggep
paradoks minangka bukti langsung yen himpunan kasebut ora
ana, banjur sinau urip kanthi kasunyatan iki. Pendekatan
iki kang bakal kita tindakake.

\end{document}
